\begin{abstract}
A world model that predicts video well may have learned something about the underlying physics, or
it may simply be good at interpolating frames. Standard probes cannot tell the difference: even
randomly initialised DreamerV3 models can contain latent directions that correlate strongly with a
pendulum's energy. We instead ask whether the model has learned a quantity that its \emph{own
dynamics} treat as conserved. In a frozen DreamerV3 trained only on $64\times64$ pendulum videos,
with no physical supervision, we recover a scalar that correlates with true energy at $0.973$ and
remains nearly constant along encoded trajectories. When the model rolls forward without new
observations, however, this quantity begins to drift.

We use three controls to test whether the recovered quantity reflects physics. First, six untrained
models span energy correlations from $0.170$ to $0.908$, showing that high correlation alone is weak
evidence; what training produces is reproducibility, with three trained models agreeing within
$0.008$. Second, when we apply the same procedure to models trained on a damped pendulum, where no
conserved energy should exist, it correctly refuses to find one: four separate statistics
distinguish conservative and dissipative models with no seed overlap. Third, the recovered quantity
is causally useful. Projecting imagined states back toward its level set reduces rollout error on
every trained seed, whereas sixty matched random polynomial constraints increase error by $53.9\%$
at the median. The recovered law outperforms every random constraint on two of three models, though
not on the third.

The internal representation is also less coordinate-specific than it first appears.
Basis-independent properties of the recovered law agree closely across independently trained models,
while selectors based on individual latent coordinates can even change sign. The same physical
structure therefore appears in different internal charts. We also report two negative results.
Requiring the recovered quantity both to stay conserved and to generate the model's flow improves
reproducibility but not accuracy, suggesting that conservation does most of the work here. And a
frequency-weighted error-accumulation mechanism that we previously established in a much smaller
recurrent model is indistinguishable from an unweighted version in DreamerV3.

Overall, a pixel-trained world model can learn a reproducible conserved quantity that its own
dynamics respect, violate during free imagination, and can be causally corrected against, while also
showing why correlation alone cannot establish that a model has learned physics.
\end{abstract}
